\section{Introduction}
\label{sec:intro}

Cartel enforcement begins before legal proof exists. An
agency deciding where to investigate rarely starts with the
communications, bid histories, leniency testimony, or
tender-level conduct evidence that would support liability.
It starts with thinner information: administrative records,
complaints, procurement anomalies, repeated interactions, and
case leads that make some firms and markets worth a closer
look. The first legal-economic problem is therefore not proof
itself. It is the allocation of proof-producing effort.

That distinction matters in procurement. The empirical
cartel-screening literature usually starts from the bid layer:
submitted bid amounts, bid ranks, within-tender dispersion,
bid histories, and sometimes bidder networks. Classic
procurement screens compare bidding patterns under competition
and collusion \citep{porter1993detection,bajari2003deciding};
variance and descriptive screens formalize suspicious bid
distributions \citep{abrantesmetz2006a,harrington2008detecting}
and field-test them in procurement settings
\citep{imhof2018screening,imhof2019detecting}. Recent work
extends the approach through bidder groups
\citep{conley2016detecting}, machine learning
\citep{huber2019machine,wallimann2023machine}, and robust
screen design \citep{chassang2022robust}.
Those objects are close to the evidentiary stage, but they are
not always cheap to obtain at scale. Many procurement systems
make contract-award records routinely visible while keeping
per-bidder bid files in operational systems that require
case-specific recovery. An enforcement agency then faces a
sequencing problem: before opening the costly bid layer, can it
use the cheap award layer to decide where bid-layer forensics
should begin?

This paper studies that question on S\~ao Paulo's BEC
procurement platform, covering $\valSampleN$ tender-items from
$\valYearStart$ to $\valYearEnd$. BEC records a routine
award layer---participants, winners, item codes, negotiated
prices, procuring units, years, and modalities---and also
preserves a richer bid layer recoverable through the LANCES
export. CADE procurement-cartel adjudications provide external
legal anchors. This institutional combination lets us study an
earlier stage than most bid-distribution screens: the point at
which an agency observes award records but has not yet recovered
the bid-level evidence needed for forensic evaluation.

The award-layer statistic is deliberately simple. Among firms
with zero wins, we rank participation intensity using
$\log(1+\text{tenders\_count})$ and implement an auditable
binary rule that flags \emph{frequent losers}. The economic
logic is monotone rather than classificatory. If coordinated
bidders use loser-side participation to create the appearance of
competition while avoiding the win, persistent zero-win
participation can rank loser-side exposure under a transparent
behavioral condition. The statistic does not identify a firm
type or assign cartel membership. In the paper's staged design,
the rank allocates forensic attention; the richer bid layer
remains where tender-level conduct and agreement are evaluated.

That legal scope determines the validation target. Direct CADE
defendants are legal membership anchors, but they are not the
natural object of a zero-win loser-side screen. Procurement
cartels often allocate winning roles, rotate designated
winners, and use other firms on the losing side of the tender
\citep{pesendorfer2000study,asker2010leniency}. The same
role-allocation logic is central in cartel theory and dynamic
case evidence \citep{marshall2012economics,clark2021collusion}.
Large-scale
procurement evidence shows how such arrangements can appear in
administrative bidding records \citep{kawai2022detecting}.
We therefore
validate the screen against always-loser firms that bid
alongside direct CADE defendants in adjudicated cartel
environments. These cobidders define an adjudication-anchored
exposure target. Direct defendants anchor legally adjudicated
environments; cobidders are the loser-side footprint the award
layer can plausibly rank before bid-level proof is recovered.

% CO-AUTHOR EDIT [v21 tone]: compress the in-intro threat catalog to a pointer;
% the full audit construction already lives in Section 4 and Appendix D.
The empirical design confronts the threats that would make such a
screen uninformative: that co-bidding concentration merely
reflects participation volume or opportunity-set exposure near
CADE defendants, that a retrospective score reuses information
from the target period, and that a loser-side statistic is
overread as a cartel-membership detector. Section~\ref{sec:validation}
meets each with participation-disciplined placebos,
exposure-adjusted audits that compare firms within more similar
opportunity sets, strictly ex ante timing exercises, and a
direct-defendant test the screen is designed to fail;
\ref{app:validation_audits_submission} reports the exact
construction.

% CO-AUTHOR EDIT [v21 tone]: lead with the result; "limited" qualifies after.
% Cobidder profile aligned with the matched evidence in Section 5: product
% concentration + bid-gap proximity survive volume matching; breadth is volume.
The frequent-loser ranking concentrates adjudication-anchored
cobidders inside the always-loser stratum, including when the
score is formed before the target window. By design it performs
poorly against direct CADE defendants---the intended boundary of
a loser-side screen, which rules out a generic cartel-membership
interpretation. Within that scope the claim is bounded but
operationally useful. Cobidders inside the frequent-loser stratum
also look different from other frequent losers: holding
participation volume fixed, they operate in more concentrated
product portfolios and bid closer to winners, while their broader
deployment is largely a volume feature. They also appear closer
to legal cartel anchors, and display bid-level patterns
consistent with, but not diagnostic of, credible losing roles.
Those patterns give the adjacency target economic content and
make it harder to reduce the result to ordinary high-volume
losing alone.

The main institutional result is a cost-of-evidence result from
the sequencing exercise. On the same adjudication-anchored
target, the award-layer score adds information to an
Imhof--Wallimann-style bid-distribution benchmark. The two
layers therefore enter at different stages: award records
prioritize where to look, and bid records evaluate what is
found. Used as a gatekeeper, the award-layer ranking cuts the
bid-microdata pool for the forensic stage by \valGateFootprintPct{} while
recovering \valGateSeqKTwoKOneTPIn{} of \valCobidders{} adjudicated cobidders. Joint scoring with
award and bid features is the full-observability upper bound;
sequential gatekeeping is the architecture relevant for agencies
that must decide where to spend costly bid recovery.

This paper makes three contributions. First, it reframes
procurement-cartel screening as evidence allocation under costly
observability. Existing bid-distribution screens show how to
evaluate suspicious bidding once bid-level data have been
recovered; we study the prior institutional stage, where
agencies observe award records but must decide where recovery
should start. This turns screening toward enforcement triage,
not toward treating a statistical flag as proof
\citep{harrington2008detecting,sanchezgraells2019screening};
robust-screen design formalizes the same distinction between
flags and evidentiary conclusions \citep{chassang2022robust}.
Second, the paper shows that a minimal award-layer
statistic---persistent zero-win participation---ranks a legally
limited loser-side adjacency target anchored in CADE
adjudications. Third, it evaluates a sequential architecture in
which cheap award-layer triage gates costly bid-layer forensics,
reducing the evidentiary footprint without moving liability out
of the richer case record.

The boundaries are design choices, not caveats. The award layer
ranks forensic priority; the validation target is
adjudication-anchored exposure; the bid layer remains the
proof-producing stage; and price evidence enters as scope
information, not a damages exercise. Because the architecture is
sequential, strategic response is part of the design: a static
cutoff invites gaming, so refreshed thresholds, continuous ranks
or capacity-constrained queues, bunching checks, and bid-layer
follow-up are its natural implementation.

The broader lesson is institutional. A legal system that waits
for full proof before ranking investigative priorities will
spend scarce forensic capacity poorly; a legal system that
treats cheap suspicion as proof will overreach. The useful
middle ground is a disciplined triage rule: one that is cheap
enough to run before proof exists, informative enough to order
forensic attention, and legally limited enough to leave
liability where it belongs---in the richer evidentiary record.
